\chapter{Filters, Matching Networks, and Transformers}
\label{ch:filters}
This chapter is the first of four that \emph{apply} the circuit models rather than
build them. Almost nothing here is derived from scratch: a filter is the RC circuit
of \cref{ch:rc} or the LC circuit of \cref{ch:rlc} with its component values chosen on
purpose, and the design rules below are statements about where those circuits' poles
and zeros sit. What is new is the \emph{intent}---given a required response, work
backwards to the components.

\section{Why Filters Matter in Amateur Radio}
Filters select desired signals and reject unwanted ones. They appear in receiver
front ends, intermediate-frequency (IF) stages, audio stages, transmitter
outputs, anti-aliasing paths,
power supplies, and interference cures. A filter is nothing more than a circuit
whose transfer function has been shaped on purpose---poles placed to form the
passband, zeros placed to notch the stopband---so everything from Part~II applies
directly.

\section{Basic Filter Types}
Four responses cover almost every filtering job, distinguished by which band of
frequencies they pass:

\begin{center}
\begin{tabularx}{\textwidth}{@{}L{0.26\textwidth}Y@{}}
\toprule
Type & What it does \\
\midrule
Low-pass & Passes low frequencies; attenuates high frequencies. \\
High-pass & Passes high frequencies; attenuates low frequencies. \\
Band-pass & Passes a band of frequencies; attenuates outside it. \\
Band-stop (notch) & Rejects a band of frequencies; passes outside it. \\
\bottomrule
\end{tabularx}
\end{center}

\section{Where Filter Behavior Comes From}
Two circuits already developed from first principles are the whole foundation of
this chapter, so we recall the results rather than re-deriving them:
\begin{description}
\item[First order --- the RC low-pass (\cref{ch:rc}).] One energy store, one pole
at \(s=-1/RC\), giving \(G(s)=1/(1+RCs)\): flat to the corner
\(\omega_c=1/RC\), then \(-\SI{20}{\decibel}\) per decade. Taken across the
resistor instead, the same circuit is a high-pass.
\item[Second order --- the LC low-pass (\cref{sec:lc-lowpass} in
\cref{ch:rlc}).] Two energy stores, a pole pair, giving
\[
G(s)=\frac{\omega_0^2}{s^2+\dfrac{\omega_0}{Q}s+\omega_0^2},
\qquad
\omega_0=\frac{1}{\sqrt{LC}},
\qquad
Q=\frac{1}{R}\sqrt{\frac{L}{C}},
\]
which falls at \(-\SI{40}{\decibel}\) per decade and peaks by \(\approx Q\) near
\(\omega_0\) when \(Q>1/\sqrt2\).
\end{description}
Those two supply everything a filter designer needs: a \emph{corner frequency}
set by the reactance values, a \emph{roll-off rate} set by the number of poles,
and---for second order and above---a \emph{peaking} set by the damping. The rest
of this chapter is about choosing them deliberately, and
\cref{sec:filterspecs} then translates the result into the numbers a data sheet
actually quotes.

\section{Order, Energy Stores, and Roll-Off}
The pattern generalizes. Each independent energy-storage element (an \(L\) or a
\(C\)) contributes one pole, so an \(n\)th-order filter has \(n\) reactive
elements and rolls off at about \(-20n\,\si{\decibel}\) per decade far above
cutoff (each pole adding \(-\SI{20}{\decibel}\)/decade, each zero
\(+\SI{20}{\decibel}\)/decade). Steeper skirts require more poles, hence more
inductors and capacitors. The \emph{order} is thus both a circuit fact (how many
reactive parts) and a response fact (how fast it rolls off).

That roll-off rule is not an approximation peculiar to filters:
\cref{sec:cascadeadd} derives it from the fact that cascaded sections multiply,
so their decibel magnitudes and their phases simply add. The same section explains
why the accompanying phase lag approaches \(-90n^\circ\)---which is what limits how
many stages can be wrapped in a feedback loop (\cref{sec:threepoles}).

\section{Zeros in Radio Circuits: Traps and Notches}
If poles build the passband, \emph{zeros} are how a circuit deliberately
\emph{rejects} energy. A transmission zero on the \(\jj\omega\) axis forces the
output to zero at one frequency. Two dual constructions put one there. A series LC
branch to ground has zero impedance at its resonance, shorting that frequency to
ground and placing a transmission zero at \(\omega_0=1/\sqrt{LC}\)---a trap. A
parallel LC branch in series with the signal path is the dual: it becomes an open
circuit at resonance, again blocking one frequency. Wherever the exam describes a
trap, a wave trap, or a notch rejecting a specific frequency, a zero has been
parked at that frequency (\cref{fig:pz-notch} in \cref{ch:splane} shows the
pole--zero picture). The elliptic family below uses exactly this---stopband
\(\jj\omega\)-axis zeros---to sharpen its skirts.

\section{Designing and Tuning a Filter}
Turn the derivation around and you have the design procedure. A filter is
specified by a target frequency response; designing it means \emph{placing the
poles and zeros} to meet that target and then realizing them with a circuit.
There are three independent choices:
\begin{description}
\item[Cutoff frequency] scale all \(L\) and \(C\) together: \(\omega_0=1/\sqrt{LC}\)
shifts the whole response along the frequency axis without changing its shape.
\item[Order] the number of poles (reactive elements) sets skirt steepness,
\(-20n\,\si{\decibel}\) per decade.
\item[Shape (the ``family'')] for a given order, \emph{where} the poles sit trades
passband flatness against transition sharpness against phase linearity.
\end{description}

That last choice is what the named families are: different rules for placing the
poles (and zeros) of a given order. \Cref{fig:filter-families} shows them at the
same order, and \cref{fig:butterworth-poles} shows the archetype---Butterworth
poles equally spaced on a circle. \Cref{sec:poleplacement} derives that pattern
and turns it into the table of section \(Q\) values a designer actually builds
from, so the families below are summarized here rather than re-derived.

\begin{tabularx}{\textwidth}{@{}L{0.19\textwidth}Y@{}}
\toprule
Family & Pole/zero placement and consequence \\
\midrule
Butterworth & Poles equally spaced on a semicircle of radius \(\omega_0\) in the left half-plane (derived in \cref{sec:poleplacement}, with the section-\(Q\) table). Maximally flat passband, no ripple, moderate transition. \\
Chebyshev & Poles on an ellipse pushed closer to the \(\jj\omega\) axis, raising every section's \(Q\) (\cref{sec:poleplacement}). Sharper cutoff at the price of passband ripple and worse phase/transient behavior. \\
Elliptic (Cauer; ``elliptical'' in the pool) & Chebyshev-like poles plus transmission \emph{zeros} on the \(\jj\omega\) axis in the stopband---depth bought at specific frequencies rather than slope everywhere. Sharpest transition; ripple in both bands. \\
Bessel & Poles positioned for maximally flat \emph{group delay} \(\tau_g=-\dd\phi/\dd\omega\), i.e.\ linear phase (defined in \cref{sec:groupdelay}); gentle amplitude cutoff. Good for pulse fidelity. \\
Crystal / cavity & Very high-\(Q\) resonators (poles hugging the axis) for narrow, highly selective band-pass shapes. \\
DSP FIR & Digital (DSP, digital signal processing) filter with no feedback (FIR, finite impulse response): no poles away from the origin, so it can be exactly linear phase---proved in \cref{sec:firiir}. \\
DSP IIR & Recursive digital filter (IIR, infinite impulse response): has poles, so it mimics analog families efficiently but with nonlinear phase (\cref{sec:firiir}). \\
\bottomrule
\end{tabularx}

\begin{figure}[htbp]
\centering
\begin{tikzpicture}[>=Latex,scale=1.15]
\draw[->,black!70] (-2.7,0) -- (0.8,0) node[right]{\(\sigma\)};
\draw[->,black!70] (0,-2.0) -- (0,2.0) node[above]{\(\jj\omega\)};
\draw[GuideLine,dashed] (0,0) circle (1.6);
% 4th-order Butterworth poles: angles 112.5,157.5,202.5,247.5 deg
\foreach \a in {112.5,157.5,202.5,247.5}{
  \filldraw[GuideBlue] ({1.6*cos(\a)},{1.6*sin(\a)}) circle (2.2pt);
}
\node[GuideLine] at (1.15,1.35) {\small \(|s|=\omega_0\)};
\node[GuideBlue] at (-2.15,1.2) {\small Butterworth};
\end{tikzpicture}
\caption{Fourth-order Butterworth poles: four left-half-plane poles equally
spaced on the circle of radius \(\omega_0\). Higher order adds more poles on the
same circle; Chebyshev pulls them toward the \(\jj\omega\) axis for a sharper but
rippled response.}
\label{fig:butterworth-poles}
\end{figure}

The \emph{transition band} is the frequency span between the passband edge and the
stopband; a ``moderate'' transition (Butterworth) is wider than the ``sharp''
transition of a Chebyshev or elliptic filter of the same order.
\Cref{fig:filter-families} makes the flatness-versus-sharpness-versus-ripple trade
concrete.

\begin{figure}[htbp]
\centering
\includegraphics[width=0.9\textwidth]{filter_families}
\caption{Order-5 low-pass responses, all normalized to the same cutoff
\(\omega_c\). Butterworth is maximally flat with a moderate transition; Chebyshev
trades \SI{1}{\decibel} of passband ripple for a sharper skirt; elliptic is
sharpest of all, with ripple in both bands and stopband nulls (its
\(\jj\omega\)-axis zeros); Bessel is gentlest, chosen for clean phase and pulse
fidelity rather than sharpness.}
\label{fig:filter-families}
\end{figure}

\begin{advancedbox}
\textbf{Where those curves come from.} Each family is defined by a formula for the
magnitude response, from which the pole positions follow. For order \(n\) and cutoff
\(\omega_c\):
\[
|G|^2_{\text{Butterworth}}=\frac{1}{1+(\omega/\omega_c)^{2n}},
\qquad
|G|^2_{\text{Chebyshev}}=\frac{1}{1+\varepsilon^2T_n^2(\omega/\omega_c)},
\]
where \(T_n\) is the order-\(n\) Chebyshev polynomial and \(\varepsilon\) sets the
ripple. \Cref{sec:poleplacement} takes the Butterworth specification apart and shows
why its poles land equally spaced on a circle of radius \(\omega_c\)
(\cref{fig:butterworth-poles}), and what section \(Q\) each conjugate pair
therefore carries. The other families are variations on that one calculation:
Chebyshev's roots lie on an ellipse; elliptic (Cauer) filters replace \(T_n\) with a
rational elliptic function, which is what introduces the finite stopband zeros; and
Bessel poles come from Bessel polynomials chosen to flatten group delay
(\cref{sec:groupdelay}). \Cref{fig:filter-families} was generated numerically from
these standard designs (via \texttt{scipy.signal}) rather than sketched, so the
ripple depths and skirt slopes are quantitatively right.
\end{advancedbox}

Finally, the same target poles can be \emph{realized} several ways: a passive LC
ladder, an active Sallen--Key stage (\cref{ch:active}) that makes complex poles
without inductors, a crystal or cavity when very high \(Q\) is needed, or a DSP
filter (FIR/IIR) that places the poles and zeros in software. What licenses that
freedom is the factoring theorem of \cref{sec:factoring}: any response of order
\(n\), however it is built, is a cascade of first- and second-order sections, so
a design reduces to a list of \((\omega_0,\zeta)\) pairs and the hardware only has
to hit them. Same poles, different hardware---so ``customizing'' a filter means choosing the order (skirt
steepness), the family (the flatness/sharpness/phase trade), the cutoff (by
scaling \(L\) and \(C\)), and the realization that suits the frequency range and
available parts.

\section{Reading a Filter's Specification}
\label{sec:filterspecs}
A data sheet does not describe a filter by its pole locations; it quotes a handful of
numbers. Each one is a measurement of the response this chapter has been shaping, and
naming them makes the datasheet and the transfer function the same object.

\begin{center}
\begin{tabularx}{\textwidth}{@{}L{0.24\textwidth}Y@{}}
\toprule
Specification & What it measures on \(|G(\jj\omega)|\) \\
\midrule
Cutoff frequency & The \(-\SI{3}{\decibel}\) (half-power) edge of the passband. For a band-pass, there are two, and their separation is the bandwidth. \\
Insertion loss & Attenuation \emph{inside} the passband---what the filter costs you on the signal you wanted. Ideally \SI{0}{\decibel}; real components make it a fraction of a decibel to a few decibels. \\
Ripple & Peak-to-peak variation within the passband: zero for Butterworth, a design choice for Chebyshev and elliptic. \\
Ultimate rejection & The attenuation floor far into the stopband. Poles give slope forever in theory, but stray coupling around the filter sets a real-world limit---often \SI{60}{\decibel} to \SI{100}{\decibel}. \\
Noise bandwidth & The width of the ideal rectangular filter passing the same noise, \(B_n=\int|H|^2\dd f/|H|^2_{\max}\). Always \emph{wider} than the \(-\SI{3}{\decibel}\) figure---by \SI{1.96}{\decibel} for one pole---because the skirts admit noise the nameplate ignores (\cref{sec:noisebw}). \\
Shape factor & The ratio of the stopband width to the passband width, conventionally \(\SI{-60}{\decibel}\) bandwidth divided by \(\SI{-6}{\decibel}\) bandwidth. It is one number for skirt steepness: \(1.0\) would be a brick wall, and smaller is sharper. \\
\bottomrule
\end{tabularx}
\end{center}

Two of these deserve comment because they are where the idealization leaks.
\emph{Insertion loss} exists because the inductors and capacitors have finite
component \(Q\) (\cref{ch:rlc})---a lossless filter would have none, and a
high-order, narrow filter built from mediocre parts can have several decibels of it.
\emph{Ultimate rejection} exists because signals get around a filter without going
through it, by stray capacitance or shared ground impedance
(\cref{sec:groundloops}); adding poles past that point buys nothing, which is why
crystal filters are physically shielded compartment by compartment.

\begin{exambox}
Shape factor and order both describe skirt steepness, and a question may give you
either. Order is the theoretical statement---\(n\) poles give
\(-20n\,\si{\decibel}\) per decade. Shape factor is the measured statement, and it
folds in the family: at the same order a Chebyshev has a smaller (better) shape
factor than a Butterworth, which is exactly the ripple-for-sharpness trade.
\end{exambox}

\section{Impedance Matching}
An impedance-matching network does two jobs:
\begin{enumerate}
\item cancel the reactive part of a complex load, and
\item transform the remaining resistance to the desired system impedance.
\end{enumerate}
Both halves are already familiar: the reactance to be cancelled is the \(X\) of
\cref{ch:ac}, and ``maximum power transfer'' is the condition derived in
\cref{ch:resistive}---\(R_L=R_{Th}\) for resistances, and the conjugate match
\(Z_L=Z_S^{*}\) once reactance is present. In a \(\SI{50}{\ohm}\) RF system, matching
maximizes power transfer and keeps transmission-line reflections controlled.

\subsection{The L-Network}
\label{sec:lnetwork}
The simplest matching network uses two reactances. To match two real resistances
\(R_{\text{hi}}>R_{\text{lo}}\), place a series reactance on the low-resistance
side and a shunt reactance on the high-resistance side. That is the whole design, and
it needs no new theory: it is \cref{sec:seriesparallel}'s series-to-parallel
transformation run backwards, once.

Hang \(X_{\text{series}}\) on \(R_{\text{lo}}\). The pair is a series branch with
\(Q=X_{\text{series}}/R_{\text{lo}}\), so by that section it is \emph{identical}, at
this one frequency, to a parallel combination of
\[
R_p=\left(1+Q^2\right)R_{\text{lo}}
\qquad\text{and}\qquad
X_p=\left(1+\frac{1}{Q^2}\right)X_{\text{series}} .
\]
Now demand that the transformed resistance be the one you wanted,
\(R_p=R_{\text{hi}}\). That single equation fixes \(Q\), and with it everything else:
\[
Q=\sqrt{\frac{R_{\text{hi}}}{R_{\text{lo}}}-1},
\qquad
X_{\text{series}}=Q\,R_{\text{lo}},
\qquad
X_{\text{shunt}}=\frac{R_{\text{hi}}}{Q}.
\]
The third follows from the second: substituting \(X_{\text{series}}=QR_{\text{lo}}\)
into \(X_p\) gives
\(X_p=R_{\text{lo}}(Q+1/Q)=R_{\text{lo}}(1+Q^2)/Q=R_{\text{hi}}/Q\). So the
transformation leaves behind a shunt reactance of exactly \(R_{\text{hi}}/Q\), and
the shunt element's only job is to cancel it---which is why the two reactances must
be of opposite type (one inductor, one capacitor). What remains is \(R_{\text{hi}}\),
purely real. \emph{An L-network is one series-to-parallel transformation plus a
cancellation}, and there is nothing else in it.

Two consequences. Choosing which element is the inductor sets whether the network is
low-pass or high-pass (\cref{fig:lnetwork}). And an L-network has no free parameter
left over---\(Q\) was determined by the resistance ratio alone---so its bandwidth is
fixed once the resistances are chosen, a restatement of
\cref{sec:seriesparallel}'s warning that the equivalence holds at one frequency.
This \(Q\) is the branch \(Q\) of the network's own series arm, the same
\(X/R\) ratio as everywhere else in the book; what makes it a \emph{network} \(Q\)
rather than a resonator's is that the \(R\) is a termination you chose, not a loss you
are stuck with.

\begin{figure}[htbp]
\centering
\begin{circuitikz}[scale=1]
% source-side termination R_lo (drawn as a resistor to ground)
\draw (0,2) coordinate(lt) to[R,l_=\(R_{\text{lo}}\)] (0,0) coordinate(lb);
% series reactance on the low-resistance side
\draw (lt) to[L,l=\(X_{\text{series}}\)] (3,2) coordinate(mid);
% shunt reactance on the high-resistance side
\draw (mid) to[C,l=\(X_{\text{shunt}}\)] (3,0) coordinate(mb);
% load-side termination R_hi (drawn as a resistor to ground)
\draw (mid) to[short] (5.5,2) coordinate(rt);
\draw (rt) to[R,l=\(R_{\text{hi}}\)] (5.5,0) coordinate(rb);
% ground rail
\draw (lb) -- (rb);
\node[ground] at (3,0){};
% input / output port labels
\node[above=0.05cm,GuideBlue] at (lt) {\small input (source side)};
\node[above=0.05cm,GuideBlue] at (rt) {\small output (load side)};
\draw[GuideBlue,-{Latex}] (-0.55,1.3) -- (-0.55,1.85);
\draw[GuideBlue,-{Latex}] (6.05,1.85) -- (6.05,1.3);
\end{circuitikz}
\caption{An L-network matching a low resistance \(R_{\text{lo}}\) to a high one
\(R_{\text{hi}}\) (both drawn as resistors to the ground rail). The network sits
\emph{between} the two terminations: the source drives the left-hand port and the
load hangs on the right-hand one. The series reactance goes on the low-resistance
side, the shunt reactance on the high-resistance side. Swapping which element is the
inductor makes it high-pass instead of low-pass.}
\label{fig:lnetwork}
\end{figure}

\subsection{Pi and T Networks}
Adding a third element gives a \(\pi\) or T network (\cref{fig:pit-network}),
which can be viewed as two L-networks back to back that meet at an internal
\emph{virtual resistance} \(R_v\).

\begin{figure}[htbp]
\centering
\begin{circuitikz}[scale=0.9]
% --- Pi network ---
\draw (0,2) coordinate(pin) to[L,l=\(L\)] (3,2) coordinate(pout);
\draw (pin) to[C,l_=\(C_1\)] (0,0) coordinate(pinb);
\draw (pout) to[C,l=\(C_2\)] (3,0) coordinate(poutb);
\draw (pinb) -- (poutb);
\node[ground] at (1.5,0){};
\draw (pin) to[short,-o] (-0.9,2) node[left]{in};
\draw (pout) to[short,-o] (3.9,2) node[right]{out};
\node at (1.5,-1.0) {\(\pi\) network};
\end{circuitikz}
\qquad
\begin{circuitikz}[scale=0.9]
% --- T network ---
\draw (0,2) coordinate(tin) to[L,l=\(L_1\)] (2,2) coordinate(tm)
      to[L,l=\(L_2\)] (4,2) coordinate(tout);
\draw (tm) to[C,l=\(C\)] (2,0) coordinate(tmb);
\draw (0,0) coordinate(tinb) -- (4,0) coordinate(toutb);
\node[ground] at (2,0){};
\draw (tin) to[short,-o] (-0.9,2) node[left]{in};
\draw (tout) to[short,-o] (4.9,2) node[right]{out};
\node at (2,-1.0) {T network};
\end{circuitikz}
\caption{Three-element matching networks. The \(\pi\) network (two shunt
capacitors, one series inductor) and the T network (two series inductors, one
shunt capacitor) each add a third element over the L-network, giving the extra
degree of freedom that sets the loaded \(Q\) through the internal virtual
resistance \(R_v\).}
\label{fig:pit-network}
\end{figure}

Choosing \(R_v\) below both terminations sets a loaded \(Q\) of roughly
\[
Q_{\text{loaded}}\approx\sqrt{\frac{R_{\max}}{R_v}-1},
\]
so \(R_v\) is the extra degree of freedom that lets the designer pick the loaded
\(Q\)---and therefore the bandwidth and harmonic rejection---independently of the
resistance ratio. ``Loaded \(Q\)'' here means exactly what it meant in
\cref{ch:rlc,ch:rlcpar}: a fractional bandwidth of roughly \(1/Q\), because that is a
statement about how far the poles sit from the \(\jj\omega\) axis and not about
wiring. This is why transmitter output and antenna-tuner networks are
usually \(\pi\) or T: they match \emph{and} set selectivity. Higher loaded \(Q\)
narrows bandwidth and improves harmonic suppression at the cost of higher
circulating currents and losses.

\begin{mistakebox}
Matching is not the same thing as filtering. A network can do both, but the
concepts are distinct: matching controls the impedance seen by a source or line,
while filtering controls frequency-dependent transmission. A \(\pi\)-network tuner
happens to be a low-pass filter too, which is a feature, not the definition.
\end{mistakebox}

\section{Transformers}
For an ideal transformer (\cref{fig:transformer}),
\[
\frac{V_s}{V_p}=\frac{N_s}{N_p},
\qquad
\frac{I_s}{I_p}=\frac{N_p}{N_s},
\qquad
Z_p=\left(\frac{N_p}{N_s}\right)^2 Z_s.
\]

\begin{figure}[htbp]
\centering
\begin{circuitikz}[scale=1]
\node[transformer core] (T) at (0,0){};
\draw (T.A1) to[short,-o] ++(-1.1,0) node[left]{\(+\)} node[left=0.7cm]{\(V_p\)};
\draw (T.A2) to[short,-o] ++(-1.1,0) node[left]{\(-\)};
\draw (T.B1) to[short,-o] ++(1.1,0) node[right]{\(+\)} node[right=0.7cm]{\(V_s\)};
\draw (T.B2) to[short,-o] ++(1.1,0) node[right]{\(-\)};
\node[above=0.15cm] at (T.A1) {\(N_p\)};
\node[above=0.15cm] at (T.B1) {\(N_s\)};
\end{circuitikz}
\caption{Ideal two-winding transformer. The primary (\(N_p\) turns) and secondary
(\(N_s\) turns) share a common core; voltages scale as the turns ratio
\(N_s/N_p\), currents inversely, and impedance as the ratio squared.}
\label{fig:transformer}
\end{figure}
These follow from flux linkage rather than from KVL alone, which is why
\cref{ch:resistive} set mutual coupling aside when it combined inductors. Impedance
transforms as the turns ratio \emph{squared}. \textbf{Core saturation}
should be avoided: every magnetic core has a maximum flux density it can carry,
and once the drive pushes it near that limit the core's permeability collapses, so
the inductance drops and the winding current shoots up. The result is distortion,
heating, and loss of the intended impedance transformation---the reason
transformers and inductors are rated for a maximum voltage-time product (or power)
at a given frequency.

\begin{exambox}
Filter questions sort into four independent decisions, and naming which one is
being asked usually settles the answer. \emph{Which band passes} gives the type
(low-pass, high-pass, band-pass, band-stop). \emph{How steep the skirt} gives the
order, at \(-20n\,\si{\decibel}\) per decade or about \(-6n\,\si{\decibel}\) per
octave. \emph{What the passband and transition look like} gives the family:
Butterworth is flat with a moderate transition, Chebyshev is sharper but rippled,
elliptic is sharpest with ripple in both bands and stopband nulls, Bessel is
gentlest but preserves waveform shape. \emph{Where a specific frequency is killed}
means a zero---a trap or notch, not a steeper filter.

Matching questions are a different animal: they are always the same two steps,
cancel the reactance and then transform the resistance. An L-network's \(Q\) is
fixed by the resistance ratio alone; a \(\pi\) or T network adds the third element
precisely so \(Q\)---and therefore bandwidth and harmonic rejection---becomes a
free choice.
\end{exambox}

\section{Worked Example: Quarter-Wave and L-Network}
\textbf{Problem.} A transmitter feeds a \(\SI{50}{\ohm}\) system, but the load
(say an antenna at resonance) presents a purely resistive \(\SI{200}{\ohm}\).
Design two matching networks that make the \(\SI{200}{\ohm}\) load look like
\(\SI{50}{\ohm}\) to the source: (a) a quarter-wave transformer, and (b) an
L-network. Compare them.

To match a real \(\SI{200}{\ohm}\) load to a \(\SI{50}{\ohm}\) system with a
lossless quarter-wave transformer (the formula \(Z_t=\sqrt{Z_0Z_L}\) is derived in
\cref{ch:lines}),
\[
Z_t=\sqrt{Z_0Z_L}=\sqrt{(50)(200)}=\SI{100}{\ohm}.
\]
The same match with an L-network gives
\[
Q=\sqrt{\frac{200}{50}-1}=\sqrt{3}\approx1.73,
\qquad
X_{\text{series}}=1.73\times50\approx\SI{87}{\ohm},
\qquad
X_{\text{shunt}}=\frac{200}{1.73}\approx\SI{116}{\ohm}.
\]
Both realize the match; the quarter-wave version needs a line one quarter
wavelength long at the design frequency, while the L-network is lumped and
narrower band.
